\section{Justificación}

La contratación pública en Colombia constituye un mecanismo fundamental para la ejecución de proyectos de infraestructura y desarrollo social. Sin embargo, la persistencia de ineficiencias y riesgos en su gestión ha generado impactos negativos tanto en el uso de los recursos públicos como en el bienestar de las comunidades.


Por lo general, los mecanismos de control y vigilancia en la contratación estatal han actuado bajo un enfoque reactivo, actuando una vez el daño patrimonial se ha realizado. Este enfoque limita la capacidad de las entidades de control para prevenir riesgos y optimizar la asignación de recursos en los procesos de auditoría. A raíz de esto, diferentes estudios han planteado la necesidad de cambiar a un esquema de control preventivo que se apoye en el análisis de datos y el uso de tecnologías emergentes, como el machine learning.

En los últimos años, el crecimiento de plataformas de datos abiertos como SECOP II ha permitido el acceso a grandes volúmenes de datos estructurados sobre procesos en la contratación pública. Estos datos incluyen información de todas las etapas en este proceso. La accesibilidad a estos datos ha impulsado el desarrollo de enfoques basados en machine learning para la detección de irregularidades, anomalías relacionadas con corrupción y problemas contractuales.  Investigaciones previas han demostrado que variables como el tipo de contratación, el número de oferentes y el valor del contrato contienen información que puede ser usado para anticipar comportamientos irregulares desde las primeras etapas del proceso de contratación \cite{gallego2021preventing} \cite{salazar2024vigia}.

Asimismo, Varios trabajos han evidenciado la efectividad de modelos de aprendizaje supervisado como Random Forest, regresión logística y XGBoost han facilitado la predicción de riesgos asociados a proyectos de infraestructura, abarcando retrasos y sobrecostos \cite{decarolis2022corruption},\cite{alshboul2022extreme},\cite{galjanic2025neural}. En el entorno colombiano, estudios que incluyen datos de plataformas públicas han identificado factores clave que se relacionan con  desviaciones de tiempo y costo, lo que resalta la viabilidad de utilizar datos abiertos para este tipo de analisis \cite{gomez2025data},\cite{gomez2025analyzing}.



Además, desde una perspectiva institucional, el uso de la inteligencia artificial en la contratación pública ha sido reconocido como una herramienta con gran capacidad para mejorar la transparencia y trazabilidad en las decisiones administrativas de igual forma reduce la irregularidades en la información que favorecen la corrupción \cite{calderon2025contratacion},\cite{sierra2025contratacion}. De este modo, los sistemas de alertas tempranas se plantean como una alternativa que puede complementar a los mecanismos de control actuales, permitiendo centrar los esfuerzos de supervisión en contratos con mayor probabilidad de riesgo, especialmente en aquellos donde los recursos son limitados. \cite{gallego2021preventing},\cite{ospina2024percepcion}.

El presente proyecto se justifica en la necesidad de integrar estos avances teóricos y tecnológicos en una propuesta aplicada al entorno del Sistema General de Regalías (SGR), el cual, concentra una proporción significativa de la inversión pública en infraestructura y ha estado históricamente asociado a problemáticas de ejecución. Pese a la disponibilidad de datos y de los avances en analitica, aun existe una brecha en el desarrollo de herramientas que permitan identificar de manera sistemática, desde la etapa de adjudicación, contratos con alto riesgo de convertirse en proyectos fallidos  .

Con esto, el desarrollo de un modelo predictivo que use técnicas de Machine Learning que sea orientado como un sistema que pueda realizar alertas tempranas, no solo representa un gran avance académico en el area de la analitica de datos aplicados a procesos estatales, sino que también una contribución práctica para fortalecer los mecanismos de control. Este tipo de herramienta puede ayudar a entidades de control y veedurías ciudadanas en la toma de decisiones,  impulsando una gestión más eficiente, transparente y responsable de los recursos públicos,en concordancia con los principios establecidos en el marco normativo de la contratación estatal \cite{ospina2024percepcion}.

Por último, este trabajo se alinea con las tendencias actuales en un gobierno abierto y el uso de datos públicos, donde la analítica avanzada se posiciona como un elemento clave para mejorar la gobernanza y la transparencia en los procesos. Como señalan varios estudios recientes, cruzar datos abiertos con Machine Learning es un avance  para el sector público. Nos permite dejar de ser reactivos y empezar a anticiparnos a los problemas antes de que ocurran \cite{lyra2022fraud}.